\documentclass[aspectratio=169]{beamer}

\usepackage[utf8]{inputenc}
\usepackage[T1]{fontenc}
\usepackage[activeacute]{babel}
\usepackage{amsmath, amssymb}
\usepackage{booktabs}
\usepackage{xcolor}
\usepackage{tcolorbox}

\usetheme{Madrid}
\usecolortheme[RGB={30,60,120}]{structure}

\setbeamertemplate{navigation symbols}{}
\setbeamertemplate{footline}{
  \leavevmode
  \hbox{
    \begin{beamercolorbox}[wd=.45\paperwidth,ht=2.5ex,dp=1ex,left,leftskip=1em]{author in head/foot}
      \usebeamerfont{author in head/foot}\insertshortauthor
    \end{beamercolorbox}
    \begin{beamercolorbox}[wd=.35\paperwidth,ht=2.5ex,dp=1ex,center]{title in head/foot}
      \usebeamerfont{title in head/foot}MAD1 -- Unidad 5
    \end{beamercolorbox}
    \begin{beamercolorbox}[wd=.2\paperwidth,ht=2.5ex,dp=1ex,right,rightskip=1em]{date in head/foot}
      \usebeamerfont{date in head/foot}\insertframenumber{} / \inserttotalframenumber
    \end{beamercolorbox}
  }
}

\definecolor{unsazul}{RGB}{30,60,120}
\definecolor{unsazulclaro}{RGB}{220,230,245}
\definecolor{verdeapp}{RGB}{20,120,60}

\tcbuselibrary{skins, breakable}

\newtcolorbox{aplicacion}{
  colback=verdeapp!8,
  colframe=verdeapp,
  fonttitle=\bfseries\small,
  title={Aplicaciones reales},
  left=4pt, right=4pt, top=4pt, bottom=4pt
}

\newtcolorbox{concepto}{
  colback=unsazulclaro,
  colframe=unsazul,
  fonttitle=\bfseries\small,
  title={Idea clave},
  left=4pt, right=4pt, top=4pt, bottom=4pt
}

\newtcolorbox{atencion}{
  colback=orange!10,
  colframe=orange!80!black,
  fonttitle=\bfseries\small,
  title={Atenci\'on},
  left=4pt, right=4pt, top=4pt, bottom=4pt
}

\title[Unidad 5 -- Representación]{Unidad 5: Aprendizaje No Supervisado}
\subtitle{Bloque 2 -- Representación de datos}
\author[MAD1 -- UNS]{Métodos de Análisis de Datos I}
\institute{Departamento de Matemática -- Universidad Nacional del Sur}
\date{}

%---------------------------------------------------------------
\begin{document}
%---------------------------------------------------------------

\begin{frame}
\titlepage
\end{frame}

%---------------------------------------------------------------
\begin{frame}{¿Por qué reducir dimensionalidad?}
\begin{columns}[T]
\column{0.55\textwidth}
Los datos modernos son de alta dimensión: imágenes ($p > 10^6$), genómica ($p > 20{,}000$), texto ($p > 100{,}000$). Esto genera problemas:

\medskip
\begin{itemize}
  \item \textbf{Maldición de la dimensionalidad:} en dimensiones altas, todos los puntos quedan ``lejos'' entre sí.
  \item \textbf{Visualización imposible:} no podemos ver más de 3 dimensiones.
  \item \textbf{Modelos inestables:} muchas variables correlacionadas generan colinealidad.
\end{itemize}

\medskip
\begin{concepto}
Los métodos de reducción buscan una función $f: \mathbb{R}^p \to \mathbb{R}^q$ con $q \ll p$ que preserve la estructura relevante de los datos.
\end{concepto}

\column{0.42\textwidth}
\textbf{Comparativa de métodos:}

\medskip
{\small
\begin{tabular}{lll}
\toprule
\textbf{Método} & \textbf{Tipo} & \textbf{Supervisado} \\
\midrule
PCA   & Lineal    & No \\
t-SNE & No lineal & No \\
UMAP  & No lineal & No \\
LDA   & Lineal    & Sí \\
\bottomrule
\end{tabular}
}

\medskip
\begin{aplicacion}
Visualización de datos genómicos, compresión de imágenes, preprocesamiento para clasificadores, exploración de corpus de texto.
\end{aplicacion}
\end{columns}
\end{frame}

%---------------------------------------------------------------
\begin{frame}{PCA: Análisis de Componentes Principales}
\begin{columns}[T]
\column{0.52\textwidth}
\textbf{Objetivo:} encontrar proyecciones lineales que maximizan la varianza.

\medskip
Sea $\mathbf{S}$ la matriz de covarianza muestral. PCA resuelve:
\[
\mathbf{v}_1 = \arg\max_{\|\mathbf{v}\|=1} \mathbf{v}^\top \mathbf{S}\, \mathbf{v}
\]
La solución es el autovector de $\mathbf{S}$ de mayor autovalor.

\medskip
\textbf{Varianza explicada} por componente $j$:
\[
\frac{\lambda_j}{\sum_{l=1}^p \lambda_l}
\]

\medskip
\begin{concepto}
PCA es como encontrar la mejor sombra de una nube de puntos 3D sobre un plano 2D: la que revela más estructura. Las componentes son esos ángulos óptimos.
\end{concepto}

\column{0.45\textwidth}
\begin{aplicacion}
\textbf{Genómica:} reducir 20.000 genes a 10--20 componentes en estudios de expresión génica. Las primeras PCs separan tipos celulares o condiciones experimentales.

\medskip
\textbf{Neurociencia:} analizar señales de cientos de electrodos simultáneos. PCA extrae los ``modos'' de activación neuronal predominantes.

\medskip
\textbf{Climatología:} reducir datos de temperatura de miles de estaciones meteorológicas. La primera PC suele capturar el calentamiento global.

\medskip
\textbf{Finanzas:} los factores de riesgo (mercado, sector, tamaño) emergen como componentes principales de los retornos de activos.

\medskip
\textbf{Imágenes:} eigenfaces -- representar rostros humanos como combinación de ``rostros base'' (PCA sobre imágenes de caras).
\end{aplicacion}
\end{columns}
\end{frame}

%---------------------------------------------------------------
\begin{frame}{PCA: selección de componentes e interpretación}
\begin{columns}[T]
\column{0.50\textwidth}
\textbf{¿Cuántas componentes retener?}

\medskip
\textbf{Varianza acumulada:} elegir $q$ tal que $\sum_{j=1}^q \lambda_j / \sum \lambda_l \geq 0.90$.

\medskip
\textbf{Scree plot:} graficar $\lambda_j$ vs.\ $j$. Cortar en el codo.

\medskip
\textbf{Regla de Kaiser:} retener solo $\lambda_j > 1$ (variables estandarizadas).

\medskip
\textbf{Interpretación:}
\begin{itemize}
  \item \textbf{Loadings} $v_{jl}$: peso de la variable $l$ en la componente $j$.
  \item \textbf{Scores} $z_{ij}$: coordenada de la observación $i$ en la componente $j$.
  \item \textbf{Biplot:} superpone scores y loadings en un gráfico.
\end{itemize}

\column{0.47\textwidth}
\begin{atencion}
PCA es sensible a la escala. \textbf{Siempre estandarizar} antes de aplicarlo, salvo que todas las variables estén en la misma unidad y se quiera preservar diferencias de varianza.
\end{atencion}

\medskip
\begin{aplicacion}
\textbf{Psicología:} en tests de inteligencia con 50 ítems, la primera PC suele ser el ``factor $g$'' (inteligencia general), que explica el 30--40\% de la varianza.

\medskip
\textbf{Química analítica:} espectros de 1.000 longitudes de onda se reducen a 5--10 PCs que representan los compuestos presentes en la muestra.

\medskip
\textbf{Deportes:} métricas de rendimiento de jugadores (pases, goles, duelos) se comprimen en 2--3 PCs que resumen el ``estilo de juego''.
\end{aplicacion}
\end{columns}
\end{frame}

%---------------------------------------------------------------
\begin{frame}{t-SNE: visualización no lineal}
\begin{columns}[T]
\column{0.52\textwidth}
\textbf{Idea:} preservar la estructura \textbf{local} de los datos en 2D.

\medskip
\begin{enumerate}
  \item Modelar similitudes en alta dimensión con gaussianas ($p_{ij}$).
  \item Modelar similitudes en 2D con distribución $t$ de Student ($q_{ij}$) --- colas pesadas para evitar colapso.
  \item Minimizar divergencia KL entre ambas distribuciones.
\end{enumerate}

\medskip
\textbf{Parámetro clave:} perplejidad (5--50), controla cuántos vecinos se consideran ``cercanos''.

\medskip
\begin{atencion}
Las \textbf{distancias entre clusters} en t-SNE \textbf{no son interpretables}: solo la estructura interna de cada cluster es confiable. Ejecutar siempre con múltiples semillas.
\end{atencion}

\column{0.45\textwidth}
\begin{aplicacion}
\textbf{Inmunología:} visualizar millones de células en experimentos de citometría de flujo. t-SNE revela subpoblaciones de linfocitos no distinguibles a simple vista.

\medskip
\textbf{NLP:} visualizar embeddings de palabras (Word2Vec, BERT). Palabras semánticamente similares forman clusters visibles en 2D.

\medskip
\textbf{Dígitos escritos a mano:} el dataset MNIST (70.000 imágenes de dígitos 0--9) forma 10 clusters bien separados en el mapa t-SNE.

\medskip
\textbf{Oncología:} visualizar perfiles de expresión génica de tumores. Cada tipo de tumor forma un cluster; los casos atípicos quedan aislados.

\medskip
\textbf{Música:} visualizar canciones por características acústicas. Géneros musicales forman regiones diferenciadas.
\end{aplicacion}
\end{columns}
\end{frame}

%---------------------------------------------------------------
\begin{frame}{UMAP: alternativa moderna a t-SNE}
\begin{columns}[T]
\column{0.52\textwidth}
UMAP (Uniform Manifold Approximation and Projection) está basado en geometría diferencial y topología algebraica.

\medskip
\textbf{Ventajas sobre t-SNE:}
\begin{itemize}
  \item Más rápido (escala a millones de puntos).
  \item Preserva mejor la \textbf{estructura global}, no solo local.
  \item Determinístico con semilla fija.
  \item Funciona para $q > 2$ dimensiones.
  \item Permite proyectar nuevas observaciones.
\end{itemize}

\medskip
\textbf{Parámetros:}
\begin{itemize}
  \item \texttt{n\_neighbors}: balance local/global (típico: 15).
  \item \texttt{min\_dist}: compacidad del embedding (típico: 0.1).
\end{itemize}

\medskip
\begin{concepto}
PCA aplana una hoja de papel. t-SNE y UMAP estiran y doblan una tela de goma para separar los grupos, aunque la distancia entre grupos no tenga significado geométrico directo.
\end{concepto}

\column{0.45\textwidth}
\begin{aplicacion}
\textbf{Biología celular:} UMAP es el estándar de facto para visualizar datos de secuenciación de ARN de célula única (scRNA-seq). Mapas de decenas de miles de células revelan trayectorias de diferenciación celular.

\medskip
\textbf{Astronomía:} el relevamiento Gaia de mil millones de estrellas usa UMAP para explorar relaciones entre parámetros estelares.

\medskip
\textbf{Ciberseguridad:} visualizar tráfico de red para detectar patrones anómalos en espacios de alta dimensión.

\medskip
\textbf{Materiales:} explorar espacios de composiciones químicas en búsqueda de nuevos materiales con propiedades deseadas.
\end{aplicacion}
\end{columns}
\end{frame}

%---------------------------------------------------------------
\begin{frame}{LDA: reducción supervisada}
\begin{columns}[T]
\column{0.52\textwidth}
LDA (Linear Discriminant Analysis) maximiza la separabilidad entre clases \textbf{conocidas}.

\medskip
\textbf{Criterio de Fisher:}
\[
J(\mathbf{w}) = \frac{\mathbf{w}^\top \mathbf{S}_B\, \mathbf{w}}{\mathbf{w}^\top \mathbf{S}_W\, \mathbf{w}}
\]

$\mathbf{S}_B$: dispersión \textbf{entre} clases.\\
$\mathbf{S}_W$: dispersión \textbf{dentro} de clases.

\medskip
La solución son los autovectores de $\mathbf{S}_W^{-1}\mathbf{S}_B$.

Con $k$ clases: hasta $\min(k-1, p)$ dimensiones.

\medskip
\begin{concepto}
PCA busca la dirección de mayor varianza total. LDA busca la dirección que mejor \textbf{separa} los grupos: que las nubes de cada clase queden juntas entre sí y lejos de las demás.
\end{concepto}

\column{0.45\textwidth}
\begin{atencion}
LDA es \textbf{supervisado}: requiere etiquetas de clase. Es útil como exploración (¿son las clases linealmente separables?) y como preprocesamiento antes de un clasificador.
\end{atencion}

\medskip
\begin{aplicacion}
\textbf{Medicina:} reducir decenas de biomarcadores a 1--2 ejes discriminantes entre ``sano'', ``pre-enfermo'' y ``enfermo''.

\medskip
\textbf{Reconocimiento facial:} Fisherfaces -- variante de LDA que supera a eigenfaces (PCA) porque maximiza la separación entre personas.

\medskip
\textbf{Enología:} clasificar vinos por origen a partir de 13 variables químicas. Con LDA, 2 ejes discriminantes separan perfectamente las tres regiones productoras.

\medskip
\textbf{Arqueometría:} separar vasijas de distintas culturas por composición química de la arcilla.

\medskip
\textbf{Marketing:} separar compradores de no compradores a partir de variables demográficas.
\end{aplicacion}
\end{columns}
\end{frame}

%---------------------------------------------------------------
\begin{frame}{Resumen: Bloque 2 -- Representación de datos}
\begin{center}
\renewcommand{\arraystretch}{1.4}
{\small
\begin{tabular}{lllll}
\toprule
\textbf{Método} & \textbf{Transformación} & \textbf{Estructura} & \textbf{Sup.} & \textbf{Uso típico} \\
\midrule
PCA   & Lineal    & Global (varianza)   & No & Preprocesamiento, exploración \\
t-SNE & No lineal & Local               & No & Visualización 2D \\
UMAP  & No lineal & Local + global      & No & Visualización, preprocesamiento \\
LDA   & Lineal    & Separabilidad       & Sí & Exploración supervisada, clasificación \\
\bottomrule
\end{tabular}
}
\end{center}

\medskip
\begin{columns}[T]
\column{0.48\textwidth}
\textbf{Regla práctica:}
\begin{itemize}
  \item Para preprocesamiento y comprensión $\to$ \textbf{PCA}.
  \item Para visualización exploratoria $\to$ \textbf{UMAP} (o t-SNE).
  \item Si tengo etiquetas y quiero ver separabilidad $\to$ \textbf{LDA}.
  \item Nunca usar t-SNE/UMAP para inferir distancias entre clusters.
\end{itemize}

\column{0.48\textwidth}
\textbf{Pipeline típico:}
\begin{enumerate}
  \item Estandarizar variables.
  \item PCA para reducir dimensión y eliminar colinealidad.
  \item UMAP para visualización del espacio reducido.
  \item LDA si hay clases y quiero ver su separabilidad.
\end{enumerate}
\end{columns}
\end{frame}

\end{document}
